\subsection{Derivation of the measured interferometric autocorrelation}
\label{sec:iac-derivation}
In interferometric autocorrelation (IAC), we measure the integrated second-harmonic signal while changing the delay between two pulse replicas. To obtain a pulse duration from this signal, we first derive the Gaussian overlap terms and then relate the optical delay to the recorded time. We use $t$ for time within a pulse, $\tau$ for the optical delay between replicas, and $x$ for the much slower oscilloscope time.

The oscilloscope records the signal as the mirror moves. The recorded-time separation $\Delta x=x_{i+1}-x_i$ is the time elapsed between two neighboring fringe maxima, expressed here in milliseconds. The optical delay between the pulses is instead on the femtosecond scale. Mirror motion connects these two times: if the optical path changes more quickly, more fringes pass in the same recorded interval. The laser's optical frequency is unchanged.

\subsubsection{Measured signal and expansion of the tutorial formula}
Assume identical replicas, an instantaneous second-order response, weak conversion, and phase matching over the relevant bandwidth. The second-harmonic field is proportional to the square of the combined fundamental field. A detector integrating over the pulse therefore records
\begin{equation}
S_{\rm SHG}(\tau)=\eta\int_{-\infty}^{\infty}
\left|[E(t)+E(t-\tau)]^2\right|^2dt
=\eta\int_{-\infty}^{\infty}|E(t)+E(t-\tau)|^4dt,
\end{equation}
where $\eta$ is a constant conversion and detection factor. This is tutorial Eq.~(4.1). For comparison, its printed expansion, Eq.~(4.2), is
\begin{equation}\tag{4.2}\label{eq:tutorial-iac}
\begin{aligned}
I_{\rm AC}(\tau)={}&\int_{-\infty}^{\infty}[I^2(t)+I^2(t-\tau)]\,dt\\
&+\int_{-\infty}^{\infty}[I(t)+I(t-\tau)]
\operatorname{Re}\{E(t)E^*(t-\tau)\}\,dt\\
&+\int_{-\infty}^{\infty}\operatorname{Re}\{E^2(t)E^{*2}(t-\tau)\}\,dt\\
&+\int_{-\infty}^{\infty}I(t)I(t-\tau)\,dt .
\end{aligned}
\end{equation}
Here $I(t)=|E(t)|^2$. We denote the four integrals by $T_0,T_1,T_2,T_3$. To find their relative coefficients, we expand the fourth power explicitly. With
\begin{equation}
a=|E(t)|^2,\qquad b=|E(t-\tau)|^2,\qquad
z=E(t)E^*(t-\tau),\qquad zz^*=ab.
\end{equation}
Then
\begin{align}
|E(t)+E(t-\tau)|^4
&=(a+b+z+z^*)^2\\
&=a^2+b^2+2ab+2(a+b)(z+z^*)+z^2+z^{*2}+2zz^*\\
&=a^2+b^2+4ab+4(a+b)\operatorname{Re}z
+2\operatorname{Re}(z^2).
\end{align}
The identities $z+z^*=2\operatorname{Re}z$ and $z^2+z^{*2}=2\operatorname{Re}(z^2)$ follow by adding conjugate pairs; no real-field assumption is made. Integration gives
\begin{equation}\label{eq:corrected-iac}
S_{\rm SHG}=\eta(T_0+4T_1+2T_2+4T_3).
\end{equation}
The expansion gives relative factors of 1, 4, 2, and 4. These factors are missing from the printed Eq.~(4.2), so we include them when calculating the signal. They cannot be absorbed into one common normalization.

\subsubsection{Evaluation of the Gaussian overlap integrals}
Take an unchirped Gaussian field with intensity standard deviation $\sigma_t$:
\begin{equation}
E(t)=E_0e^{-t^2/(4\sigma_t^2)}e^{-i\omega_0t},
\qquad I(t)=I_0e^{-t^2/(2\sigma_t^2)},\qquad I_0=|E_0|^2.
\end{equation}
The common Gaussian integral can be evaluated explicitly. Since its integrand is positive,
\begin{align}
J&=\int_{-\infty}^{\infty}e^{-v^2}\,dv>0,\\
J^2&=\int_{\mathbb R^2}e^{-(v^2+w^2)}\,dv\,dw
=\int_0^{2\pi}d\theta\int_0^\infty e^{-\rho^2}\rho\,d\rho\\
&=2\pi\left[-\tfrac12e^{-\rho^2}\right]_0^\infty=\pi,
\qquad J=\sqrt{\pi}.
\end{align}
For any real shift $d$, substituting $v=(t-d)/\sigma_t$, so $dt=\sigma_t\,dv$, gives
\begin{equation}
\int_{-\infty}^{\infty}e^{-(t-d)^2/\sigma_t^2}dt
=\sigma_t\int_{-\infty}^{\infty}e^{-v^2}dv
=\sqrt{\pi}\sigma_t.
\end{equation}
Define $Q=I_0^2\sqrt{\pi}\sigma_t$. Translation invariance gives the constant term:
\begin{equation}
T_0=I_0^2\int_{-\infty}^{\infty}
\left[e^{-t^2/\sigma_t^2}+e^{-(t-\tau)^2/\sigma_t^2}\right]dt=2Q.
\end{equation}
For the fundamental-frequency term, the carrier phases multiply as
\begin{equation}
z=E(t)E^*(t-\tau)
=I_0e^{-[t^2+(t-\tau)^2]/(4\sigma_t^2)}
e^{-i\omega_0t}e^{i\omega_0(t-\tau)}
=I_0e^{-[t^2+(t-\tau)^2]/(4\sigma_t^2)}e^{-i\omega_0\tau}.
\end{equation}
In the definition of $T_1$, the real part is taken before integration. It can be moved outside the integral because $w(t)=I(t)+I(t-\tau)$ is real and the Gaussian integrals converge absolutely. At fixed $\tau$,
\begin{align*}
T_1&=\int_{-\infty}^{\infty}w(t)\operatorname{Re}z(t)\,dt\\
&=\frac12\int_{-\infty}^{\infty}w(t)[z(t)+z^*(t)]\,dt\\
&=\frac12\left[\int_{-\infty}^{\infty}w(t)z(t)\,dt
+\left(\int_{-\infty}^{\infty}w(t)z(t)\,dt\right)^*\right]\\
&=\operatorname{Re}\int_{-\infty}^{\infty}w(t)z(t)\,dt.
\end{align*}
The third line uses $\int w z^*\,dt=(\int w z\,dt)^*$ for real $w$. Thus the complex overlaps may be integrated first and their real parts taken afterward, without assuming that $z$ is real.

Completing the square separately for its two intensity factors gives
\begin{align}
3t^2+(t-\tau)^2&=4(t-\tau/4)^2+3\tau^2/4,\\
t^2+3(t-\tau)^2&=4(t-3\tau/4)^2+3\tau^2/4.
\end{align}
Consequently,
\begin{align}
\int_{-\infty}^{\infty}I(t)z\,dt
&=I_0^2e^{-i\omega_0\tau}e^{-3\tau^2/(16\sigma_t^2)}
\int_{-\infty}^{\infty}e^{-(t-\tau/4)^2/\sigma_t^2}dt\\
&=Qe^{-3\tau^2/(16\sigma_t^2)}e^{-i\omega_0\tau},\\
\int_{-\infty}^{\infty}I(t-\tau)z\,dt
&=I_0^2e^{-i\omega_0\tau}e^{-3\tau^2/(16\sigma_t^2)}
\int_{-\infty}^{\infty}e^{-(t-3\tau/4)^2/\sigma_t^2}dt\\
&=Qe^{-3\tau^2/(16\sigma_t^2)}e^{-i\omega_0\tau}.
\end{align}
Taking the real part of their sum, using $\operatorname{Re}e^{-i\theta}=\cos\theta$, yields
\begin{equation}
T_1=2Qe^{-3\tau^2/(16\sigma_t^2)}\cos(\omega_0\tau).
\end{equation}
For the remaining two terms, use
\begin{equation}
t^2+(t-\tau)^2=2(t-\tau/2)^2+\tau^2/2.
\end{equation}
The intensity overlap becomes
\begin{align}
T_3&=I_0^2\int_{-\infty}^{\infty}
e^{-[t^2+(t-\tau)^2]/(2\sigma_t^2)}dt\\
&=I_0^2e^{-\tau^2/(4\sigma_t^2)}
\int_{-\infty}^{\infty}e^{-(t-\tau/2)^2/\sigma_t^2}dt
=Qe^{-\tau^2/(4\sigma_t^2)}.
\end{align}
Since $z^2=I(t)I(t-\tau)e^{-2i\omega_0\tau}$, the second-harmonic overlap is
\begin{equation}
T_2=\operatorname{Re}\int_{-\infty}^{\infty}z^2dt
=\operatorname{Re}\{e^{-2i\omega_0\tau}T_3\}
=Qe^{-\tau^2/(4\sigma_t^2)}\cos(2\omega_0\tau).
\end{equation}
Inserting all four terms into Eq.~\eqref{eq:corrected-iac} gives
\begin{equation}\label{eq:gaussian-iac}
\begin{aligned}
S_{\rm SHG}(\tau)=\eta Q\bigl\{2
&+8e^{-3\tau^2/(16\sigma_t^2)}\cos(\omega_0\tau)\\
&+e^{-\tau^2/(4\sigma_t^2)}
[4+2\cos(2\omega_0\tau)]\bigr\}.
\end{aligned}
\end{equation}
The cosine terms follow directly from the relative carrier phase of the two replicas. At zero delay, the result is $16\eta Q$, while for separated pulses it tends to $2\eta Q$. This gives the expected ideal interferometric contrast of $8:1$.

\subsubsection{Pulse FWHM versus autocorrelation-envelope width}
The pulse duration quoted in this report is the \emph{intensity} FWHM. Its half-maximum condition gives
\begin{align}
I(t_{1/2})/I_0=e^{-t_{1/2}^2/(2\sigma_t^2)}&=\tfrac12,\\
t_{1/2}^2&=2\sigma_t^2\ln2,\\
\Delta t_{\rm pulse}=2|t_{1/2}|&=2\sqrt{2\ln2}\,\sigma_t.
\end{align}
For the two distinct Gaussian overlap factors in Eq.~\eqref{eq:gaussian-iac},
\begin{align}
g(\tau)&=e^{-\tau^2/(4\sigma_t^2)},&
h(\tau)&=e^{-3\tau^2/(16\sigma_t^2)},\\
g(\tau_{1/2})=\tfrac12
&\Rightarrow |\tau_{1/2}|=2\sigma_t\sqrt{\ln2},&
\Delta\tau_g&=4\sigma_t\sqrt{\ln2}
=\sqrt2\,\Delta t_{\rm pulse},\\
h(\tau_{1/2})=\tfrac12
&\Rightarrow |\tau_{1/2}|=4\sigma_t\sqrt{\ln2/3},&
\Delta\tau_h&=8\sigma_t\sqrt{\ln2/3}
=\sqrt{8/3}\,\Delta t_{\rm pulse}.
\end{align}
The two overlap functions therefore have different widths. The factor $\sqrt2$ relates the isolated intensity autocorrelation $T_3=Qg$ to the pulse duration. Since the full interferometric signal contains both $g$ and $h$, applying this factor to the width of its complete envelope would give an incorrect duration.

The tutorial's alternative field notation $E(t)\propto e^{-t^2/\Gamma^2}$ corresponds to $\Gamma=2\sigma_t$. Hence
\begin{equation}
\Delta t_{\rm pulse}=\Gamma\sqrt{2\ln2},
\qquad
\Delta t_{\rm field}=2\Gamma\sqrt{\ln2}.
\end{equation}
Thus $\Gamma$ itself is neither of these FWHM values. The spectral-intensity bandwidth and its time--bandwidth product are treated separately in the preceding subsection.

\subsubsection{Recorded-time calibration and the fitted signal}
The recorded time $x$ is expressed in milliseconds relative to an initial peak estimate. Let $x_0$ be the fitted center and $r=x-x_0$. The optical-delay coordinate $u$, in carrier cycles, is modeled locally by
\begin{align}
u(r)&=fr+b_2r^2+b_3r^3,&
\tau(x)&=\frac{\lambda_0}{c}u(r),\\
\frac{du}{dx}&=f+2b_2r+3b_3r^2,&
\omega_0\tau(x)&=2\pi u(r).
\end{align}
Here $f$, $b_2$, and $b_3$ have units of cycles/ms, cycles/ms$^2$, and cycles/ms$^3$, respectively. One fundamental fringe cycle corresponds to a delay change $\lambda_0/c$. A constant scan speed would give $b_2=b_3=0$; the changing fringe spacing in Fig.~\ref{fig:fringe-calibration} motivates allowing a nonlinear mapping. Successive maxima initialize it, and the fitted derivative is checked to be positive throughout the sampled fit window. This is an empirical calibration rather than an independent measurement of mirror displacement.

The blue points in Fig.~\ref{fig:fringe-calibration} are calculated from pairs of successive detected maxima in the raw signal. For maxima at recorded times $x_i$ and $x_{i+1}$, their horizontal coordinate and vertical value are
\begin{equation}
x_{{\rm mid},i}=\frac{x_i+x_{i+1}}{2},
\qquad f_i=\frac{1}{x_{i+1}-x_i}.
\end{equation}
Thus each point represents the inverse fringe spacing, in cycles per millisecond, placed halfway between its two peaks. Closer maxima give a higher value. The code uses SciPy's \texttt{find\_peaks} function within the fitted window, with a prominence threshold of 5\% of that window's raw-signal range and a minimum separation of approximately $0.6/f$, converted to sample indices. Here $f$ is the fitted central fringe rate. The peak positions come from the measured samples, but their selection is therefore not fully independent of the fit. Noise, missed maxima, or extra detected maxima can affect these estimates; pointwise uncertainties have not been quantified. The red curve instead shows $du/dx$ from the fitted cubic mapping. Agreement supports the delay calibration, but does not independently establish a particular mechanical trajectory.

Write the pulse-intensity standard deviation in optical cycles as $s=c\sigma_t/\lambda_0$. The overlap functions then become
\begin{equation}
g(u)=e^{-u^2/(4s^2)},\qquad
h(u)=e^{-3u^2/(16s^2)}.
\end{equation}
The implemented fit uses every raw sample within \SI{2.8}{ms} of the initial peak estimate:
\begin{equation}\label{eq:fitted-iac}
\begin{aligned}
S_{\rm fit}(x)={}&B+A_0g(u)
+h(u)[C_1\cos(2\pi u)+D_1\sin(2\pi u)]\\
&+g(u)[C_2\cos(4\pi u)+D_2\sin(4\pi u)].
\end{aligned}
\end{equation}
The baseline and harmonic coefficients are independent fit parameters. In the ideal model, with matching signal units and phase origin, they would satisfy
\begin{equation}
B=2\eta Q,\quad A_0=4\eta Q,\quad C_1=8\eta Q,\quad
C_2=2\eta Q,\quad D_1=D_2=0.
\end{equation}
In the fit, we allow these coefficients to vary so that the measured amplitudes and phases need not match the ideal values. The Gaussian overlap widths are retained. This makes the fit more flexible, but does not determine the detector response or account fully for chirp.

The fitted width is converted directly to the pulse duration:
\begin{equation}
\sigma_t=s\lambda_0/c,\qquad
\boxed{\Delta t_{\rm pulse}=2\sqrt{2\ln2}\,s\lambda_0/c.}
\end{equation}
We thus obtain the pulse duration directly from the fitted $s$. The conversion does not require measuring an envelope width. The associated statistical errors and the remaining calibration and model uncertainties are discussed in Sec.~\ref{sec:uncertainty}.
